%% CHAPITRE 5 -- Securite integree : SAST, SCA, DAST, Zero Trust et IAM

\section*{Introduction}

Ce chapitre constitue le cœur technique du projet en termes de sécurité. Il couvre l'ensemble
des couches de protection mises en place~: les tests de sécurité statiques et dynamiques
intégrés dans le pipeline CI/CD (SAST, SCA, DAST), l'architecture réseau Zero Trust avec les
NetworkPolicies Kubernetes, la surveillance runtime via Wazuh SIEM, et la gouvernance des
identités et des accès (IAM). Pour chaque mécanisme, les résultats concrets obtenus sont
présentés avec les métriques mesurables.

\section{Taxonomie des Tests de Securite}
Trois types complementaires de tests de securite sont integres dans la plateforme~:
\begin{itemize}
  \item \textbf{SAST} (\textit{Static Application Security Testing})~: analyse le code source
        sans l'executer. Outils~: SonarCloud, Trivy FS, ESLint-security~;
  \item \textbf{SCA} (\textit{Software Composition Analysis})~: analyse les dependances tierces
        (packages npm, couches OS des images Docker). Outil~: Trivy~;
  \item \textbf{DAST} (\textit{Dynamic Application Security Testing})~: teste l'application en
        cours d'execution en boite noire. Outil~: OWASP ZAP.
\end{itemize}
\section{SAST avec SonarCloud}
SonarCloud analyse le code TypeScript/React a chaque push et detecte~:
\begin{itemize}
  \item \textbf{Security Hotspots}~: utilisations de \texttt{dangerouslySetInnerHTML} (XSS)~;
  \item \textbf{Vulnerabilities}~: patterns de code pouvant mener a des injections~;
  \item \textbf{Code Smells}~: indicateurs de qualite masquant des vulnerabilites~;
  \item \textbf{Coverage gate}~: bloque la PR si la couverture de tests chute sous le seuil.
\end{itemize}
\section{SCA avec Trivy --- Vulnerabilites Detectees}
Le projet integre intentionnellement des dependances vulnerables~:
\begin{table}[h]
\centering
\setlength{\arrayrulewidth}{0.8pt}
\setlength{\dashlinedash}{3pt}
\setlength{\dashlinegap}{2pt}
\renewcommand{\arraystretch}{1.4}
\caption{CVEs detectees par Trivy FS sur les dependances npm}
\label{tab:cves_npm}
\begin{tabularx}{\textwidth}{|>{\centering\arraybackslash}X|>{\centering\arraybackslash}X|>{\centering\arraybackslash}X|>{\centering\arraybackslash}X|}
\hline
\tableheader \textbf{CVE} & \textbf{Package} & \textbf{Severite} & \textbf{Description} \\
\hline
CVE-2021-23337 & lodash 4.17.20 & HIGH & Command Injection \\ \hdashline
CVE-2026-4800  & lodash 4.17.20 & HIGH & Code Injection via \_.template \\ \hdashline
CVE-2025-13465 & lodash 4.17.20 & MEDIUM & Prototype Pollution \\ \hdashline
CVE-2022-23539 & jsonwebtoken 8.5.1 & HIGH & Insecure key types \\ \hdashline
CVE-2022-23540 & jsonwebtoken 8.5.1 & MEDIUM & Signature bypass (none algo) \\ \hdashline
CVE-2022-23541 & jsonwebtoken 8.5.1 & MEDIUM & RSA to HMAC forgery \\ \hdashline
CVE-2024-29041 & express 4.17.1 & MEDIUM & Open Redirect \\
\hline
\end{tabularx}
\end{table}
Le scan de l'image \texttt{nginx:1.21} revele \textbf{47 vulnerabilites}~: 3 CRITICAL, 12 HIGH,
18 MEDIUM, 14 LOW. Les plus critiques incluent CVE-2022-0778 (OpenSSL infinite loop) et
CVE-2021-3711 (OpenSSL buffer overflow).
\section{DAST avec OWASP ZAP}
OWASP ZAP effectue un scan en boite noire de l'application deployee apres chaque deploiement~:
\begin{table}[h]
\centering
\setlength{\arrayrulewidth}{0.8pt}
\setlength{\dashlinedash}{3pt}
\setlength{\dashlinegap}{2pt}
\renewcommand{\arraystretch}{1.4}
\caption{Findings OWASP ZAP sur l'application}
\label{tab:zap_findings}
\begin{tabularx}{\textwidth}{|>{\centering\arraybackslash}X|>{\centering\arraybackslash}X|>{\centering\arraybackslash}X|}
\hline
\tableheader \textbf{Finding} & \textbf{Regle} & \textbf{Severite} \\
\hline
Stored XSS (dangerouslySetInnerHTML) & 40012 & HIGH \\ \hdashline
Missing Content-Security-Policy & 10038 & MEDIUM \\ \hdashline
Missing X-Frame-Options (Clickjacking) & 10020 & MEDIUM \\ \hdashline
Directory Browsing Enabled & 10033 & MEDIUM \\ \hdashline
Server Version Leaked (nginx/1.21.6) & 10036 & LOW \\ \hdashline
Missing HSTS & 10035 & LOW \\ \hdashline
Missing X-Content-Type-Options & 10021 & LOW \\
\hline
\end{tabularx}
\end{table}
Apres durcissement de \texttt{nginx.conf} (ajout des headers de securite, desactivation de
\texttt{server\_tokens} et \texttt{autoindex})~: \textbf{0 FAIL, 0 WARN}.
\section{Reseau Zero Trust --- NetworkPolicies}
\subsection{Principe Default-Deny}
Sans NetworkPolicies, tous les pods d'un cluster Kubernetes peuvent communiquer librement.
Ce comportement par defaut est contraire au principe du moindre privilege. La politique
appliquee~: \textbf{tout est interdit par defaut, chaque flux autorise est declare explicitement}.
Une politique \texttt{default-deny-all} est appliquee dans chaque namespace (\texttt{app},
\texttt{testing}, \texttt{observability}, \texttt{security}).
\subsection{Matrice des Flux Autorises}
\begin{table}[h]
\centering
\setlength{\arrayrulewidth}{0.8pt}
\setlength{\dashlinedash}{3pt}
\setlength{\dashlinegap}{2pt}
\renewcommand{\arraystretch}{1.4}
\caption{Matrice des flux reseau autorises}
\label{tab:network_matrix}
\begin{tabularx}{\textwidth}{|>{\centering\arraybackslash}X|>{\centering\arraybackslash}X|>{\centering\arraybackslash}X|>{\centering\arraybackslash}X|}
\hline
\tableheader \textbf{Source} & \textbf{Destination} & \textbf{Port} & \textbf{Politique} \\
\hline
\texttt{testing} pods & \texttt{app} pods & TCP 80 & allow-testing-to-app \\ \hdashline
PostSync hook & \texttt{app} pods & TCP 80 & allow-playwright-hook \\ \hdashline
\texttt{observability} & \texttt{app} sidecar & TCP 9113 & allow-prometheus-scrape \\ \hdashline
\texttt{security} DaemonSet & Wazuh VM & TCP 1514/1515 & allow-wazuh-egress \\ \hdashline
\texttt{app} pods & kube-dns & UDP/TCP 53 & allow-dns-egress \\ \hdashline
kubelet & \texttt{app} pods & TCP (probes) & allow-kubelet-probes \\
\hline
\end{tabularx}
\end{table}
\section{Securite Runtime --- Wazuh SIEM}
\subsection{Architecture Wazuh}
Wazuh est un SIEM open-source qui fournit~: detection d'intrusion (HIDS), monitoring de
l'integrite des fichiers (FIM), conformite reglementaire (PCI DSS, GDPR, NIST), et
correlation d'evenements en temps reel.
L'agent Wazuh est deploye comme \textbf{DaemonSet} dans le namespace \texttt{security},
garantissant une instance par noeud. Il collecte les logs depuis \texttt{/var/log/containers/}
et les envoie via TCP 1514 vers le manager sur la VM GCE.
\subsection{Limites GKE Autopilot}
GKE Autopilot bloque \texttt{hostPID}, \texttt{hostNetwork} et les conteneurs privilegies,
ce qui limite les capacites de monitoring kernel de Wazuh (pas d'audit syscall, pas de visibilite
reseau host). Les logs conteneurs et l'inventaire reseau/OS restent fonctionnels.
\section{Gouvernance IAM et Identite}
\subsection{Probleme des Cles JSON Statiques}
Les cles JSON GCP n'ont pas de date d'expiration par defaut, leur rotation est manuelle
et souvent negligee, et leur compromission donne un acces longue duree a l'infrastructure.
\subsection{Workload Identity Federation}
Le flux WIF en 5 etapes~:
\begin{enumerate}
  \item GitHub Actions obtient un token OIDC (JWT signe par GitHub)~;
  \item Le JWT est presente au pool WIF GCP~;
  \item GCP verifie l'issuer, le repository et la branche~;
  \item GCP emet des credentials temporaires (valides 1 heure maximum)~;
  \item GitHub Actions utilise ces credentials pour appeler les APIs GCP.
\end{enumerate}
\subsection{Modele IAM Complet}
\begin{table}[h]
\centering
\setlength{\arrayrulewidth}{0.8pt}
\setlength{\dashlinedash}{3pt}
\setlength{\dashlinegap}{2pt}
\renewcommand{\arraystretch}{1.4}
\caption{Identites et roles IAM du projet}
\label{tab:iam}
\begin{tabularx}{\textwidth}{|>{\centering\arraybackslash}X|>{\centering\arraybackslash}X|>{\centering\arraybackslash}X|>{\centering\arraybackslash}X|}
\hline
\tableheader \textbf{Identite} & \textbf{Type} & \textbf{Roles GCP} & \textbf{Usage} \\
\hline
gke-app-sa & GSA & artifactregistry.reader & Pull d'images depuis les noeuds \\ \hdashline
github-actions-sa & GSA & container.developer, artifactregistry.writer & Pipeline CI/CD \\ \hdashline
playwright-runner & KSA & Via Workload Identity & Jobs Playwright en cluster \\
\hline
\end{tabularx}
\end{table}

\section*{Conclusion}

Ce chapitre a démontré la mise en œuvre concrète des couches de sécurité~: neuf CVEs détectées
par Trivy sur les dépendances npm, 47 vulnérabilités identifiées dans l'image nginx, dix
findings reportés par OWASP ZAP, un réseau Zero Trust avec douze NetworkPolicies, un SIEM
Wazuh opérationnel, et une gouvernance IAM keyless. L'ensemble constitue une couverture
sécuritaire multi-niveaux et auditable. Le chapitre final présente les tests fonctionnels,
l'observabilité, et la synthèse quantitative des résultats obtenus.

